\documentclass[aps,pra,onecolumn,amsmath,amssymb,superscriptaddress,longbibliography]{revtex4-2}

\usepackage{booktabs}
\usepackage{bm}
\usepackage{hyperref}

\newcommand{\ii}{\mathrm{i}}
\newcommand{\ee}{\mathrm{e}}
\newcommand{\per}{\operatorname{per}}
\newcommand{\evec}{\bm e}

\begin{document}

\title{Supplemental Material for\\
``Directional response jets of suppressed multiphoton events\\
in a four-mode Fourier interferometer''}

\author{Hainan Zhao}
\email{hainzhao@gmail.com}
\affiliation{Independent Researcher}

\date{\today}
\maketitle

\section{Phase histograms and generating coefficients}
\label{supp:hist}

Let $q^4=1$ and
\begin{equation}
L_j(q)=\sum_{k=0}^{3}q^{jk}x_k .
\end{equation}
For input $\bm r$ and output $\bm s$, define
\begin{equation}
\Phi_q(\bm r,\bm s)=
\left(\prod_{k=0}^{3}r_k!\right)
[\bm x^{\bm r}]\prod_{j=0}^{3}L_j(q)^{s_j}.
\label{supp:eq:phi}
\end{equation}
If $h_\ell(\bm r,\bm s)$ counts labelled permanent paths with phase
$\ii^\ell$, then
\begin{equation}
\Phi_q(\bm r,\bm s)=\sum_{\ell=0}^{3}h_\ell(\bm r,\bm s)q^\ell .
\end{equation}
The four evaluations at $1,\ii,-1,-\ii$ determine the histogram by Fourier
inversion.

\section{Full proof of sectorwise reciprocity}
\label{supp:reciprocity}

Fix nonnegative $d,N$, let $0\leq a,s\leq d$, and let
$0\leq p,k\leq N$.  We compare
\begin{align}
\bm r&=(a,k,d-a,N-k),&
\bm s&=(s,p,d-s,N-p),\nonumber\\
\widetilde{\bm r}&=(a,p,d-a,N-p),&
\widetilde{\bm s}&=(s,k,d-s,N-k).
\label{supp:eq:pairs}
\end{align}

At $q=\ii$, set
\begin{equation}
E=x_0+x_2,\quad T=x_0-x_2,\quad
S=x_1+x_3,\quad D=x_1-x_3.
\end{equation}
The output generating polynomial factors as
\begin{align}
A(E,S)&=(E+S)^s(E-S)^{d-s}
=\sum_{m=0}^{d}\alpha_m E^{d-m}S^m,\nonumber\\
B_p(T,D)&=(T+\ii D)^p(T-\ii D)^{N-p}.
\end{align}
Use the unnormalized binary Krawtchouk coefficients
\begin{equation}
K_m(p;N)=
\sum_t(-1)^t\binom pt\binom{N-p}{m-t}.
\end{equation}
Direct binomial expansion gives
\begin{align}
[T^m]B_p
&=(-1)^{N-p-m}\ii^{N-m}K_m(p;N)D^{N-m},
\label{supp:eq:three-coeffs-a}\\
[x_0^ax_2^{d-a}]E^{d-m}T^m
&=(-1)^mK_a(m;d),
\label{supp:eq:three-coeffs-b}\\
[x_1^kx_3^{N-k}]S^mD^{N-m}
&=(-1)^{N-m-k}K_k(m;N).
\label{supp:eq:three-coeffs-c}
\end{align}
Let $C^{(a)}_{p,k}$ denote the coefficient before the input-factorial
multiplier.  Equations
\eqref{supp:eq:three-coeffs-a}--\eqref{supp:eq:three-coeffs-c} yield
\begin{multline}
C^{(a)}_{p,k}
=\sum_m\alpha_m
(-1)^{N-p-m}\ii^{N-m}
(-1)^mK_a(m;d)\\
\times(-1)^{N-m-k}
K_m(p;N)K_k(m;N).
\label{supp:eq:coefficient-sum}
\end{multline}
All factors other than the final two Krawtchouk coefficients are unchanged
under $p\leftrightarrow k$.  The standard duality
\begin{equation}
\binom NpK_m(p;N)=\binom NmK_p(m;N)
\label{supp:eq:duality}
\end{equation}
therefore gives, term by term,
\begin{equation}
C^{(a)}_{p,k}
=\frac{\binom Nk}{\binom Np}C^{(a)}_{k,p}.
\label{supp:eq:ratio}
\end{equation}
The input-factorial multipliers in Eq.~\eqref{supp:eq:phi} are
\begin{equation}
a!(d-a)!k!(N-k)!
=\frac{a!(d-a)!N!}{\binom Nk}
\end{equation}
and the same expression with $p$ replacing $k$.  They cancel the ratio in
Eq.~\eqref{supp:eq:ratio}.

At $q=-1$, the output polynomial is
\begin{equation}
(E+S)^d(E-S)^N,
\end{equation}
independent of $p$.  The requested monomials have total degree $d$ in the
even-indexed variables $(x_0,x_2)$ and total degree $N$ in the odd-indexed
variables $(x_1,x_3)$, so only the $E^dS^N$ term contributes.  Its two
odd-mode coefficients have ratio $\binom Nk/\binom Np$, canceled by the
same factorials.  At $q=1$, both labelled values count all $n!$ permanent
paths.  The $q=-\ii$ equality is the complex conjugate of the $q=\ii$
equality.  Fourier inversion proves equality of all four phase bins.

Ordinary input--output exchange, which holds because the canonical $F_4$ is
symmetric, supplies the other two equalities in the fourfold statement.
The product of input and output factorials is the same for the reciprocal
pairs, so normalized Fock amplitudes are also equal.

\subsection{Polynomial extension}

For
\begin{equation}
H(z)=
\begin{pmatrix}
1&1&1&1\\
1&z&-1&-z\\
1&-1&1&-1\\
1&-z&-1&z
\end{pmatrix},
\end{equation}
the four generating forms are $E+S,T+zD,E-S,T-zD$.  Replacing $\ii$ by
$z$ in Eq.~\eqref{supp:eq:three-coeffs-a} contributes the common factor
$z^{N-m}$ and leaves Eq.~\eqref{supp:eq:ratio} unchanged.  Reciprocal
permanents are thus identical polynomials in $z$.  The passive-unitary
interpretation is restricted to $|z|=1$.

\subsection{Boundary of the result}

The histogram is defined in the canonical dephased gauge; arbitrary row and
column phases need not preserve the four bins.  Nor does parity alone imply
an analogous theorem for $F_8$.  For example, the $F_8$ phase histograms for
\begin{align}
(0,0,1,0,0,1,0,0)&\longrightarrow(1,0,0,1,0,0,0,0),\nonumber\\
(0,0,1,1,0,0,0,0)&\longrightarrow(1,0,0,0,0,1,0,0)
\end{align}
are unequal, although the odd-mode occupation subvectors are exchanged in
the naive sense.  The two binary sectors exposed by the order-four
factorization, rather than parity alone, are essential.

\section{Directional derivatives and Krylov moments}
\label{supp:moments}

For the unnormalized root matrix $2F_4$, write
\begin{equation}
Z_{\bm r,\bm t}=\per(2F_4)[\bm t,\bm r],
\qquad
D_{\bm r,\bm s}=2^n
\sqrt{\prod_jr_j!\prod_ks_k!}.
\end{equation}
Second quantization maps $X_{pq}$ and $Y_{pq}$ to
\begin{align}
\widehat X_{pq}&=b_p^\dagger b_q+b_q^\dagger b_p,\nonumber\\
\widehat Y_{pq}&=-\ii b_p^\dagger b_q+\ii b_q^\dagger b_p.
\end{align}
For a dark target, direct use of the ladder factors gives
\begin{align}
\mathcal A_X'(0)
&=\frac{\ii}{D_{\bm r,\bm s}}
\left(s_qZ_{\bm r,\bm s+\evec_p-\evec_q}
+s_pZ_{\bm r,\bm s-\evec_p+\evec_q}\right),
\label{supp:eq:x-tangent}\\
\mathcal A_Y'(0)
&=\frac{1}{D_{\bm r,\bm s}}
\left(s_pZ_{\bm r,\bm s-\evec_p+\evec_q}
-s_qZ_{\bm r,\bm s+\evec_p-\evec_q}\right).
\label{supp:eq:y-tangent}
\end{align}
The square roots cancel against the normalization of the neighboring Fock
state.  At higher order, move one particle at a time in the output
occupation graph.  A move out of a mode with current occupation $m$
contributes $m$; for $Y$, it also contributes $+\ii$ or $-\ii$ according
to the orientation.  Summing the resulting Gaussian-integer weights times
$Z_{\bm r,\bm t}$ gives the common-denominator numerator of
$\langle\bm s|\widehat G^jF_4|\bm r\rangle$.

The two-mode sector reached from a target with $M=s_p+s_q$ has dimension
$M+1$.  If the first $M+1$ Krylov moments vanish, Cayley--Hamilton implies
that the amplitude vanishes for the entire one-parameter rotation.

\section{Complete event-level response tables}

An entry $c\epsilon^{2j}$ means
$P(\epsilon)=c\epsilon^{2j}+O(\epsilon^{2j+1})$.

\subsection{Cyclic four-photon event}

\begin{equation}
\bm r=(1,1,1,1),\qquad \bm s=(3,1,0,0).
\end{equation}
\begin{center}
\begin{tabular}{ccc}
\toprule
pair&$X_{pq}$&$Y_{pq}$\\
\midrule
01&$3\epsilon^2/8$&$3\epsilon^2/8$\\
02&exact&exact\\
03&$3\epsilon^2/8$&$3\epsilon^2/8$\\
12&exact&exact\\
13&exact&exact\\
23&exact&exact\\
\bottomrule
\end{tabular}
\end{center}

\subsection{Reciprocity-closed four-photon event}

\begin{equation}
\bm r=\bm s=(0,1,2,1).
\end{equation}
\begin{center}
\begin{tabular}{ccc}
\toprule
pair&$X_{pq}$&$Y_{pq}$\\
\midrule
01&$\epsilon^2/64$&$\epsilon^2/64$\\
02&$\epsilon^2/4$&$\epsilon^2/4$\\
03&$\epsilon^2/64$&$\epsilon^2/64$\\
12&$\epsilon^2/64$&$25\epsilon^2/64$\\
13&$\epsilon^2/4$&exact\\
23&$\epsilon^2/64$&$25\epsilon^2/64$\\
\bottomrule
\end{tabular}
\end{center}

For pair $13$, direct evaluation at arbitrary angle gives the unnormalized
coefficient $4\ii\sin(2\epsilon)$.  Fock normalization therefore yields
\begin{equation}
\mathcal A_{X_{13}}(\epsilon)=
\frac{\ii\sin(2\epsilon)}{4},\qquad
\mathcal A_{Y_{13}}(\epsilon)=0.
\end{equation}

\subsection{Eleven-photon affine-family root}

\begin{equation}
\bm r=(0,1,3,7),\qquad \bm s=(1,3,3,4).
\end{equation}
\begin{center}
\begin{tabular}{ccc}
\toprule
pair&$X_{pq}$&$Y_{pq}$\\
\midrule
01&$595\epsilon^2/8192$&$455\epsilon^2/16384$\\
02&$315\epsilon^2/16384$&$315\epsilon^2/8192$\\
03&$875\epsilon^2/16384$&$315\epsilon^2/16384$\\
12&$315\epsilon^2/16384$&$315\epsilon^4/8192$\\
13&$315\epsilon^2/16384$&$315\epsilon^2/16384$\\
23&$35\epsilon^2/4096$&$35\epsilon^2/4096$\\
\bottomrule
\end{tabular}
\end{center}
For $Y_{12}$, the two first-neighbor permanents are equal:
\begin{equation}
Z_{\bm r,(1,2,4,4)}
=Z_{\bm r,(1,4,2,4)}=241920.
\end{equation}
Since $s_1=s_2=3$, Eq.~\eqref{supp:eq:y-tangent} vanishes.  The exact
second-order common-denominator moment is
\begin{equation}
2903040(1-\ii).
\end{equation}
Including the Taylor factor $1/2!$ and Fock normalization gives
$315\epsilon^4/8192$.

All three traceless diagonal $SU(4)$ directions preserve every tabled zero:
an appended diagonal unitary multiplies the target amplitude by
$\exp(\ii\sum_js_j\theta_j)$.

\section{Proof of the all-odd exact $Y_{13}$ axis}
\label{supp:odd-proof}

Let $\bm r=\bm s=(0,a,2a,a)$, with $a$ odd.  Put
$x=x_1$, $y=x_2$, $z=x_3$, $u=x+z$, and $v=x-z$.
For $c=\cos\epsilon$, $d=\sin\epsilon$, define
$\alpha=c+d$ and $\beta=c-d$.  The two mixed Fourier row forms obey
\begin{equation}
L'_1L'_3=
\alpha\beta(y^2+v^2)
+\ii(\alpha^2-\beta^2)yv,
\qquad L_2=y-u.
\end{equation}
The desired coefficient is proportional to
\begin{equation}
[x^ay^{2a}z^a]
\left(A(y^2+v^2)+Byv\right)^a(y-u)^{2a},
\label{supp:eq:odd-coefficient}
\end{equation}
where $A=\alpha\beta$ and $B=\ii(\alpha^2-\beta^2)$.

The functional $[x^az^a]$ is invariant under $x\leftrightarrow z$, which
sends $v\mapsto-v$.  It kills every term with an odd number of $Byv$
factors.  Choose $2h$ such factors.  If $k-h$ of the remaining factors
contribute $v^2$, the central coefficient is
\begin{equation}
K_{a,k}=[x^az^a](x-z)^{2k}(x+z)^{2a-2k},
\end{equation}
with weight proportional to
\begin{equation}
\binom{a-2h}{k-h}\binom{2a}{2k}.
\end{equation}
The weight is invariant under $k\leftrightarrow a-k$.  Substituting
$z\mapsto-z$ gives
\begin{equation}
K_{a,a-k}=(-1)^aK_{a,k}.
\end{equation}
For odd $a$, each $k$ contribution cancels its $a-k$ partner.  Hence
Eq.~\eqref{supp:eq:odd-coefficient} vanishes for every $\epsilon$.

\section{Nuisance-floor and finite-shot calculation}
\label{supp:nuisance}

The illustrative model in the main text is
\begin{equation}
P_{\rm obs}=\mathcal V P_{\rm ind}
+(1-\mathcal V)P_{\rm dist}+B.
\end{equation}
For mutually distinguishable labelled photons, $P_{\rm dist}$ is the
coefficient of $\bm y^{\bm s}$ in
\begin{equation}
\prod_{k=0}^{3}
\left(\sum_{j=0}^{3}|U_{jk}|^2y_j\right)^{r_k}.
\label{supp:eq:dist}
\end{equation}
At $\epsilon=0.1$ for the four-photon $13$ comparison, the exact program
gives
\begin{align}
P_{\rm ind}^{X}&=0.00246684393741,&
P_{\rm ind}^{Y}&=0,\nonumber\\
P_{\rm dist}^{X}&=0.046875,&
P_{\rm dist}^{Y}&=0.0474917109844.
\end{align}
Substitution of $\mathcal V=0.99$ and $B=0$ gives the values in the main
text.  For $M$ independent accepted trials at each setting, the difference
of empirical proportions has approximate variance
\begin{equation}
\frac{P_X(1-P_X)+P_Y(1-P_Y)}{M}.
\end{equation}
Equating the mean contrast to $Z$ standard deviations gives the reported
normal-approximation estimate.  Exact binomial likelihoods should replace
this estimate when counts are very small.

Uniform or mode-dependent attenuation placed strictly after the complete
interferometer sends a target amplitude to
\begin{equation}
\mathcal A_{\bm r,\bm s}\prod_j\eta_j^{s_j/2}
\end{equation}
and cannot lift a zero when the full photon number is postselected.
Loss interleaved with the interferometer is different: it can change the
relative weights of interfering paths.  The present model does not claim to
identify that loss, general partial distinguishability, source impurity, or
detector backgrounds from a single observed floor.

\section{Reproducibility}

The repository uses only the Python standard library.  The relevant commands
are
\begin{verbatim}
python3 -m unittest discover -s tests -v
python3 scripts/analyze_reciprocity_census.py
python3 scripts/analyze_unitary_leakage.py
python3 scripts/analyze_finite_shot_protocol.py
\end{verbatim}
The phase histograms and leakage moments are exact integer or rational
calculations.  Floating-point arithmetic is used only for the illustrative
finite-angle nuisance model; the script separately certifies the exact
$\sin^2(2\epsilon)/16$ law to numerical tolerance at multiple angles.

\end{document}
